% ===== PRÉAMBULE CANONIQUE — à coller VERBATIM en tête de CHAQUE .tex =====
\documentclass[11pt,a4paper]{article}
\usepackage[utf8]{inputenc}\usepackage[T1]{fontenc}\usepackage{lmodern}
\usepackage[french]{babel}
\usepackage{amsmath,amssymb,mathtools}\usepackage{icomma}
\usepackage[version=4]{mhchem}          % \ce{...} : équations & formules
\usepackage{siunitx}\sisetup{locale=FR} % \qty{}{}, \si{}, \ang{}
\usepackage{chemfig}                    % \chemfig{...} : structures développées
\usepackage{xcolor}\usepackage{colortbl}
\usepackage{tikz}\usepackage{pgfplots}\pgfplotsset{compat=1.18}
\usepackage[most]{tcolorbox}\usepackage{enumitem}\usepackage{array,booktabs}
\usepackage[a4paper,margin=2.2cm]{geometry}
\usetikzlibrary{arrows.meta,calc,angles,quotes,positioning,patterns,backgrounds,shapes.geometric}
\definecolor{primary}{RGB}{222,49,99}\definecolor{primarydark}{RGB}{150,22,55}
\definecolor{defcol}{RGB}{0,114,178}\definecolor{methcol}{RGB}{52,92,148}
\definecolor{excol}{RGB}{0,135,90}\definecolor{attncol}{RGB}{200,40,55}
\tcbset{boxbase/.style={enhanced,breakable,boxrule=0.9pt,arc=2.5pt,
  left=3mm,right=3mm,top=2.3mm,bottom=2.3mm,fonttitle=\bfseries,coltitle=white}}
% --- boîtes TUTORIEL (noms EXACTS, ne pas renommer) ---
\newtcolorbox{cle}[1][]{boxbase,colback=defcol!6,colframe=defcol,
  title={\textsc{À retenir}\ifx\relax#1\relax\relax\else\textmd{ : \itshape #1}\fi}}
\newtcolorbox{methode}[1][]{boxbase,colback=methcol!7,colframe=methcol,sharp corners,
  title={\textsc{Méthode}\ifx\relax#1\relax\relax\else\textmd{ --- \itshape #1}\fi}}
\newtcolorbox{exemplebox}[1][]{boxbase,colback=excol!5,colframe=excol,
  title={\textsc{Exemple}\ifx\relax#1\relax\relax\else\textmd{ : \itshape #1}\fi}}
\newtcolorbox{piege}[1][]{boxbase,colback=attncol!6,colframe=attncol,
  title={\textsc{Piège}\ifx\relax#1\relax\relax\else\textmd{ : \itshape #1}\fi}}
\newtcolorbox{remarque}[1][]{boxbase,colback=black!4,colframe=black!45,
  title={\textsc{Remarque}\ifx\relax#1\relax\relax\else\textmd{ : \itshape #1}\fi}}
% --- boîtes EXERCICE / CORRIGÉ (noms EXACTS, auto-comptées) ---
\newtcolorbox[auto counter]{exo}[1][]{boxbase,colback=primary!4,colframe=primary,
  title={\textsc{Exercice}~\thetcbcounter\ifx\relax#1\relax\relax\else\textmd{ --- \itshape #1}\fi}}
\newtcolorbox[auto counter]{corrige}[1][]{boxbase,colback=excol!4,colframe=excol,
  title={\textsc{Corrigé}~\thetcbcounter\ifx\relax#1\relax\relax\else\textmd{ --- \itshape #1}\fi}}
\newcommand{\R}{\mathbb{R}}\newcommand{\N}{\mathbb{N}}\newcommand{\Z}{\mathbb{Z}}\newcommand{\ds}{\displaystyle}
% (un \usepackage{fancyhdr}+en-tête est possible ; il est ignoré côté web)
% =========================================================================
\begin{document}
\begin{exo}[toutes les charges de \ce{CO3^2-}]
Dans l'ion carbonate, le carbone central est lié à trois oxygènes : \emph{un} par une double
liaison, \emph{deux} par des liaisons simples (chaque \ce{O} simplement lié porte $3$ doublets
libres, l'\ce{O} doublement lié en porte $2$).
\begin{enumerate}[label=\alph*)]
  \item Charge formelle de l'oxygène doublement lié.
  \item Charge formelle d'un oxygène simplement lié.
  \item Charge formelle du carbone.
  \item Vérifie que la somme redonne la charge de l'ion.
\end{enumerate}
\end{exo}

\begin{exo}[toutes les charges de \ce{NO3-}]
Structure analogue : azote central lié à $3$ oxygènes (une double, deux simples).
\begin{enumerate}[label=\alph*)]
  \item Charge formelle de l'azote ($0$ doublet libre, $4$ liaisons).
  \item Charge formelle de l'oxygène doublement lié.
  \item Charge formelle d'un oxygène simplement lié.
  \item Somme des charges formelles.
\end{enumerate}
\end{exo}

\begin{exo}[une neutre à charges cachées : \ce{CO}]
Le monoxyde de carbone s'écrit \ce{C#O} (triple liaison, un doublet libre sur chaque atome).
\begin{enumerate}[label=\alph*)]
  \item Charge formelle du carbone.
  \item Charge formelle de l'oxygène.
  \item Que vaut la somme ? Est-ce cohérent avec le fait que \ce{CO} est neutre ?
\end{enumerate}
\end{exo}

\begin{exo}[l'ozone \ce{O3}]
L'ozone s'écrit \ce{O=O-O} : l'oxygène central porte $1$ doublet libre, l'oxygène doublement lié
$2$ doublets, l'oxygène simplement lié $3$ doublets.
\begin{enumerate}[label=\alph*)]
  \item Charge formelle de l'oxygène central.
  \item Charge formelle de l'oxygène doublement lié.
  \item Charge formelle de l'oxygène simplement lié.
  \item Somme des charges (l'ozone est neutre).
\end{enumerate}
\end{exo}

\begin{exo}[choisir la meilleure structure : \ce{CO2}]
On hésite entre deux structures de \ce{CO2} :
\begin{center}
\textbf{(i)} \ce{O=C=O} \qquad et \qquad \textbf{(ii)} \ce{O#C-O}
\end{center}
\begin{enumerate}[label=\alph*)]
  \item Donne les charges formelles des trois atomes dans la structure (i).
  \item Donne les charges formelles des trois atomes dans la structure (ii). \\
        (Dans (ii) : l'\ce{O} triplement lié a $1$ doublet libre ; l'\ce{O} simplement lié en a $3$.)
  \item Laquelle est la meilleure structure, et pourquoi ?
\end{enumerate}
\end{exo}

\end{document}
